\documentclass[11pt,a4paper]{article}
\usepackage[utf8]{inputenc}
\usepackage[T2A]{fontenc}
\usepackage[russian]{babel}
\usepackage{amsmath,amssymb,amsthm}
\usepackage{booktabs}
\usepackage{geometry}
\usepackage{hyperref}
\usepackage{mathtools}
\usepackage{tikz}
\usepackage{float}

\usetikzlibrary{shapes.geometric, arrows.meta, positioning}
\usetikzlibrary{calc}

\geometry{margin=2.5cm}

\newtheorem{definition}{Определение}
\newtheorem{remark}{Замечание}

\title{Математические основы и параметры CRYSTALS-Kyber}
\author{}
\date{}

\begin{document}

\maketitle

% \newpage
\section{Введение}

CRYSTALS-Kyber (стандартизован NIST как ML-KEM) — это механизм инкапсуляции ключа (KEM), стойкость которого основывается на сложности задачи \emph{Module Learning with Errors} (MLWE) над модульными решётками. Схема предлагает три уровня стойкости: Kyber512, Kyber768 и Kyber1024, которые ориентированы на уровни безопасности NIST 1, 3 и 5 соответственно (аналогично AES-128, AES-192 и AES-256).

Сначала Kyber строится как IND-CPA-стойкая схема шифрования с открытым ключом (Kyber.CPAPKE), а затем с помощью преобразования Фуджисаки-Окамото (FO) превращается в IND-CCA-стойкий KEM (Kyber.CCAKEM).

\newpage
\section{Математический аппарат}

\subsection{Кольца и распределения}

Все вычисления производятся в кольце
\[
    R_q = \mathbb{Z}_q[X]/(X^{256}+1),
\]
где \(q = 3329\) — простой модуль, одинаковый для всех версий Kyber. Элементы кольца — многочлены степени меньше 256 с коэффициентами из \(\mathbb{Z}_q\).

Секретные и шумовые многочлены выбираются из \emph{центрированного биномиального распределения} \(B_\eta\), определяемого как
\[
    B_\eta = \sum_{i=1}^{\eta} (a_i - b_i),
\]
где \(a_i, b_i\) независимо и равномерно выбираются из \(\{0,1\}\). Параметр \(\eta\) контролирует размах шума: коэффициенты лежат в диапазоне \([-\eta, \eta]\).

\subsection{Задача MLWE}

Стойкость Kyber опирается на сложность задачи Module-LWE. В общем виде для модуля ранга \(k\) она формулируется так: даны матрица \(\mathbf{A} \in R_q^{k \times k}\) и вектор \(\mathbf{b} = \mathbf{A}\mathbf{s} + \mathbf{e}\), где \(\mathbf{s}, \mathbf{e}\) — векторы из \(R_q^k\) с малыми коэффициентами (выбранными из \(B_\eta\)), требуется отличить \(\mathbf{b}\) от равномерного случайного элемента \(R_q^k\). Именно сложность этой задачи гарантирует, что по открытому ключу невозможно восстановить секрет \(\mathbf{s}\).

\newpage
\section{Схема Kyber.CPAPKE}

Базовая IND-CPA-схема состоит из трёх алгоритмов.

\subsection{Генерация ключей}

Выбираются случайные \(\mathbf{s}, \mathbf{e} \leftarrow B_\eta^k\). Открытый ключ вычисляется как
\[
    \mathbf{t} = \mathbf{A}\mathbf{s} + \mathbf{e} \in R_q^k,
\]
а секретным ключом является \(\mathbf{s}\). Публикуется пара \((\mathbf{A}, \mathbf{t})\), которая представляет собой экземпляр задачи MLWE.

\subsection{Шифрование (Encapsulation в CPAPKE)}

Для шифрования сообщения \(m \in \{0,1\}^{256}\) (кодируемого как многочлен с коэффициентами \(0\) или \(\lfloor q/2 \rfloor\)) выбираются эфемерные шумовые векторы
\[
    \mathbf{r}, \mathbf{e}_1 \leftarrow B_{\eta_1}^k, \qquad e_2 \leftarrow B_{\eta_2}.
\]
Шифртекст состоит из двух компонент:
\begin{align}
    \mathbf{u} & = \mathbf{A}^T \mathbf{r} + \mathbf{e}_1 \in R_q^k, \label{eq:u}                                       \\
    v          & = \mathbf{t}^T \mathbf{r} + e_2 + \left\lfloor \frac{q}{2} \right\rfloor \cdot m \in R_q. \label{eq:v}
\end{align}
Здесь \(\mathbf{e}_1\) — вектор шумовых многочленов размерности \(k\), а \(e_2\) — одиночный шумовой многочлен. Пара \((\mathbf{u}, v)\) и есть шифртекст.

\subsection{Расшифровка}

Получатель, зная секрет \(\mathbf{s}\), вычисляет
\[
    v - \mathbf{s}^T \mathbf{u}
    = \mathbf{t}^T \mathbf{r} + e_2 + \left\lfloor \frac{q}{2} \right\rfloor m - \mathbf{s}^T (\mathbf{A}^T \mathbf{r} + \mathbf{e}_1).
\]
Подставляя \(\mathbf{t} = \mathbf{A}\mathbf{s} + \mathbf{e}\), получаем
\begin{align*}
    v - \mathbf{s}^T \mathbf{u}
     & = (\mathbf{A}\mathbf{s} + \mathbf{e})^T \mathbf{r} + e_2 + \left\lfloor \frac{q}{2} \right\rfloor m
    - \mathbf{s}^T \mathbf{A}^T \mathbf{r} - \mathbf{s}^T \mathbf{e}_1                                                                             \\
     & = \mathbf{s}^T \mathbf{A}^T \mathbf{r} + \mathbf{e}^T \mathbf{r} + e_2 + \left\lfloor \frac{q}{2} \right\rfloor m
    - \mathbf{s}^T \mathbf{A}^T \mathbf{r} - \mathbf{s}^T \mathbf{e}_1                                                                             \\
     & = \left\lfloor \frac{q}{2} \right\rfloor m + \underbrace{\mathbf{e}^T \mathbf{r} - \mathbf{s}^T \mathbf{e}_1 + e_2}_{\text{шум } \Delta m}.
\end{align*}
Ключевое наблюдение: все слагаемые в \(\Delta m\) являются произведениями и суммами многочленов с коэффициентами из \(B_\eta\), поэтому коэффициенты \(\Delta m\) малы. Если они не превосходят \(\lfloor q/4 \rfloor\), то округление \(v - \mathbf{s}^T \mathbf{u}\) к ближайшему кратному \(q/2\) восстанавливает \(m\) без ошибки.

\newpage
\section{Роль \(\mathbf{e}_1\) и \(e_2\)}

Вектор \(\mathbf{e}_1\) и скаляр \(e_2\) — это \emph{шумовые добавки}, которые маскируют алгебраическую структуру шифртекста. Без них компоненты \(\mathbf{u}\) и \(v\) были бы точными линейными комбинациями \(\mathbf{r}\) и \(\mathbf{t}\), что позволило бы атакующему легко восстановить сообщение или секрет. Добавление \(\mathbf{e}_1\) и \(e_2\) превращает задачу восстановления \(\mathbf{s}\) или \(m\) в задачу о нахождении малого решения в системе линейных уравнений над \(R_q\) — то есть в вариант задачи MLWE/LWE.

Именно величины \(\eta_1\) и \(\eta_2\), задающие размах \(B_{\eta_1}\) и \(B_{\eta_2}\), определяют компромисс между стойкостью и вероятностью ошибки расшифровки. Чем больше шум, тем сложнее атака, но тем выше риск того, что \(\Delta m\) превысит порог и расшифровка окажется неверной.

\newpage
\section{Преобразование в IND-CCA KEM (Kyber.CCAKEM)}

Схема CPAPKE сама по себе не является CCA-стойкой. Kyber использует модифицированное преобразование FO: при инкапсуляции случайность \(\mathbf{r}\) детерминированно выводится из сообщения \(m\) и хэша открытого ключа, а при декапсуляции выполняется повторное шифрование и сравнение с полученным шифртекстом. Если они не совпадают, вместо расшифрованного ключа возвращается псевдослучайное значение \(z\) (implicit rejection). Это предотвращает атаки с адаптивно выбранными шифртекстами.

\newpage
\section{Параметры версий Kyber}

Все три версии используют одно и то же кольцо \(R_q\) с \(n = 256\) и \(q = 3329\). Различия заключаются в ранге модуля \(k\), размахе шума \(\eta_1, \eta_2\) и параметрах сжатия \(d_{\mathbf{u}}, d_v\).

\begin{table}[H]
    \centering
    \caption{Параметры Kyber для трёх уровней стойкости}
    \begin{tabular}{lccccccc}
        \toprule
        Версия    & Уровень NIST & \(k\) & \(\eta_1\) & \(\eta_2\) & \(d_{\mathbf{u}}\) & \(d_v\) & Вероятность ошибки \\
        \midrule
        Kyber512  & 1            & 2     & 3          & 2          & 10                 & 4       & \(2^{-139}\)       \\
        Kyber768  & 3            & 3     & 2          & 2          & 10                 & 4       & \(2^{-165}\)       \\
        Kyber1024 & 5            & 4     & 2          & 2          & 11                 & 5       & \(2^{-175}\)       \\
        \bottomrule
    \end{tabular}
\end{table}

\begin{itemize}
    \item \textbf{\(k\)} — ранг модуля. Определяет размерность векторов \(\mathbf{s}, \mathbf{e}, \mathbf{r}, \mathbf{e}_1\) и, следовательно, размер открытого ключа и шифртекста. Большее \(k\) означает более высокую стойкость, но и большие накладные расходы.
    \item \textbf{\(\eta_1, \eta_2\)} — параметры центрированного биномиального распределения. Для Kyber512 используется \(\eta_1=3\), что даёт чуть больший шум в эфемерных векторах \(\mathbf{r}, \mathbf{e}_1\); для Kyber768 и Kyber1024 \(\eta_1=2\). Параметр \(\eta_2=2\) одинаков для всех версий и задаёт шум \(e_2\).
    \item \textbf{\(d_{\mathbf{u}}, d_v\)} — количество битов, до которого сжимаются компоненты шифртекста \(\mathbf{u}\) и \(v\). Сжатие уменьшает размер шифртекста, но вносит дополнительный шум (ошибку сжатия), который учитывается при выборе параметров.
\end{itemize}

% \newpage
% \section{Во что упирается стойкость}

% Стойкость Kyber \emph{сводится} к сложности задачи MLWE: если существует алгоритм, эффективно решающий MLWE для данного кольца и распределения шума, то он может быть использован для взлома Kyber. На сегодняшний день лучшие классические и квантовые атаки на Kyber не эксплуатируют модульную структуру и рассматривают задачу как обычную LWE-задачу над решёткой размерности \(k \cdot 256\).

% Конкретные уровни стойкости (в битах) для трёх версий приведены ниже.

% \begin{table}[H]
%     \centering
%     \caption{Оценки классической и квантовой стойкости (в битах)}
%     \begin{tabular}{lccc}
%         \toprule
%         Версия    & Уровень NIST & Классическая стойкость & Квантовая стойкость \\
%         \midrule
%         Kyber512  & 1            & 118                    & 107                 \\
%         Kyber768  & 3            & 183                    & 166                 \\
%         Kyber1024 & 5            & 256                    & 232                 \\
%         \bottomrule
%     \end{tabular}
% \end{table}

% Эти цифры отражают сложность решения задачи MLWE соответствующими алгоритмами решёточного анализа (например, BKZ). Важно понимать, что стойкость упирается не в конкретную алгебраическую атаку, а в \emph{отсутствие} эффективных алгоритмов для задачи MLWE при данных параметрах \( (n=256, q=3329, k, \eta) \). Именно поэтому выбор \(k\) и \(\eta\) является центральным при настройке безопасности: увеличение \(k\) экспоненциально усложняет атаку, а увеличение \(\eta\) делает шум менее предсказуемым, но одновременно повышает вероятность ошибки расшифровки, которая должна оставаться пренебрежимо малой (порядка \(2^{-139}\) и ниже).

% \section{Заключение}

% Кратко математическую суть Kyber можно выразить так: безопасность KEM эквивалентна сложности восстановления малого секрета \(\mathbf{s}\) из пары \((\mathbf{A}, \mathbf{A}\mathbf{s}+\mathbf{e})\). Шумовые векторы \(\mathbf{e}_1\) и \(e_2\), добавляемые при шифровании, гарантируют, что даже легитимный получатель всегда имеет дело с зашумлёнными линейными комбинациями, а параметры \(k, \eta_1, \eta_2\) подобраны так, чтобы, с одной стороны, сделать задачу неразрешимой для известных атак, а с другой — удержать вероятность ошибки расшифровки на уровне, не создающем практических проблем.


\newpage
\section{Распределения коэффициентов шумовых векторов}

В Kyber все шумовые полиномы выбираются из центрированного биномиального распределения \(B_\eta\), которое определяется как
\[
    B_\eta = \sum_{i=1}^{\eta} (a_i - b_i),
\]
где \(a_i,b_i \in \{0,1\}\) — независимые равномерные биты. Область значений: \(\{-\eta,-\eta+1,\dots,0,\dots,\eta-1,\eta\}\), среднее равно \(0\), дисперсия \(\eta/2\).

Параметр \(\eta_1\) используется для векторов \(\mathbf{e}\) (шум при генерации ключей) и \(\mathbf{r}\) (случайный вектор при шифровании). Параметр \(\eta_2\) используется для векторов \(\mathbf{e}_1\) и скаляра \(e_2\) (шумы при шифровании).

\subsection{Значения \(\eta_1\) и \(\eta_2\) по уровням}

\begin{table}[H]
    \centering
    \begin{tabular}{lcc}
        \toprule
        Уровень    & \(\eta_1\) & \(\eta_2\) \\
        \midrule
        Kyber-512  & 3          & 2          \\
        Kyber-768  & 2          & 2          \\
        Kyber-1024 & 2          & 2          \\
        \bottomrule
    \end{tabular}
    \caption{Параметры распределений для трёх уровней Kyber.}
\end{table}

\subsubsection{Таблица вероятностей для \(\eta=2\)}

Для \(\eta=2\) каждый коэффициент принимает значения из \(\{-2,-1,0,1,2\}\) с вероятностями, приведёнными ниже. Это распределение используется для векторов \(\mathbf{e}_1\) и \(e_2\) на \textbf{всех} уровнях, а также для \(\mathbf{e}\) и \(\mathbf{r}\) на уровнях Kyber-768 и Kyber-1024.

\begin{table}[H]
    \centering
    \begin{tabular}{c c}
        \toprule
        Значение коэффициента & Вероятность \\
        \midrule
        \(-2\)                & \(1/16\)    \\
        \(-1\)                & \(4/16\)    \\
        \(0\)                 & \(6/16\)    \\
        \(1\)                 & \(4/16\)    \\
        \(2\)                 & \(1/16\)    \\
        \bottomrule
    \end{tabular}
    \caption{Распределение \(B_2\) (значения и вероятности).}
\end{table}

\subsubsection{Таблица вероятностей для \(\eta=3\)}

Для \(\eta=3\) область значений расширяется до \(\{-3,-2,-1,0,1,2,3\}\). Это распределение используется для векторов \(\mathbf{e}\) и \(\mathbf{r}\) \textbf{только} на уровне Kyber-512.

\begin{table}[H]
    \centering
    \begin{tabular}{c c}
        \toprule
        Значение коэффициента & Вероятность \\
        \midrule
        \(-3\)                & \(1/64\)    \\
        \(-2\)                & \(6/64\)    \\
        \(-1\)                & \(15/64\)   \\
        \(0\)                 & \(20/64\)   \\
        \(1\)                 & \(15/64\)   \\
        \(2\)                 & \(6/64\)    \\
        \(3\)                 & \(1/64\)    \\
        \bottomrule
    \end{tabular}
    \caption{Распределение \(B_3\) (значения и вероятности).}
\end{table}

\subsection{Сводная таблица: какой вектор какое распределение использует}

\begin{table}[H]
    \centering
    \begin{tabular}{lcccc}
        \toprule
        Уровень    & \(\mathbf{e}\) (вектор) & \(\mathbf{r}\) (вектор) & \(\mathbf{e}_1\) (вектор) & \(e_2\) (скаляр) \\
        \midrule
        Kyber-512  & \(B_3\)                 & \(B_3\)                 & \(B_2\)                   & \(B_2\)          \\
        Kyber-768  & \(B_2\)                 & \(B_2\)                 & \(B_2\)                   & \(B_2\)          \\
        Kyber-1024 & \(B_2\)                 & \(B_2\)                 & \(B_2\)                   & \(B_2\)          \\
        \bottomrule
    \end{tabular}
    \caption{Распределения для каждого шумового элемента по уровням.}
\end{table}

\newpage
\section{Особенности центрированного биномиального распределения и мотивация его выбора}

\subsection{Почему в Kyber отказались от гауссова распределения}

В ранних схемах на основе задачи LWE (Learning With Errors) шум и секрет обычно выбирались из дискретного гауссова распределения. Такой выбор удобен для теоретических доказательств безопасности: для LWE с гауссовым шумом существуют строгие сведения о сложности (worst-case to average-case reduction). Однако у гауссова распределения есть два серьёзных практических недостатка.

Во-первых, \textbf{выборка из гауссова распределения вычислительно неэффективна}. Она требует либо вычислений с высокой точностью (например, обратной функции ошибок), либо хранения больших таблиц предвычисленных значений. И то, и другое плохо ложится на ограниченные ресурсы встраиваемых устройств, на которые во многом ориентированы постквантовые KEM.

Во-вторых, \textbf{гауссова выборка уязвима к атакам по времени (timing attacks)}. Время выполнения операции может зависеть от выбираемого значения, что даёт атакующему побочный канал. Хотя существуют техники защиты (например, выборка с постоянным временем), они усложняют реализацию и снижают производительность.

Именно поэтому в Kyber (и в ряде других современных схем, таких как Saber) было решено заменить гауссово распределение на \textbf{центрированное биномиальное распределение (CBD, centered binomial distribution)}.

\subsection{Ключевая особенность: безопасность LWE зависит от стандартного отклонения, а не от точной формы распределения}

Важное теоретическое наблюдение, обосновывающее допустимость такой замены, состоит в следующем: \textbf{стойкость LWE-шифрования определяется не точной формой распределения шума, а его стандартным отклонением} (и, шире, моментом второго порядка). Это означает, что если CBD имеет то же (или достаточно близкое) стандартное отклонение, что и «эталонное» гауссово распределение, то замена не приводит к потере гарантий безопасности, вытекающих из LWE.

Для CBD с параметром \(\eta\) стандартное отклонение равно
\[
    \sigma_{\mathrm{CBD}} = \sqrt{\frac{\eta}{2}}.
\]
Например, при \(\eta=2\) получаем \(\sigma \approx 1{,}0\), при \(\eta=3\) — \(\sigma \approx 1{,}225\). Эти значения находятся в том же диапазоне, который использовался бы и для гауссова шума в аналогичных параметризациях LWE.

\subsection{Преимущества CBD с точки зрения реализации}

Центрированное биномиальное распределение \(B_\eta\) определяется как сумма \(\eta\) независимых разностей двух равномерных битов:
\[
    B_\eta = \sum_{i=1}^{\eta} (a_i - b_i), \qquad a_i,b_i \leftarrow \{0,1\}.
\]
Из этого определения вытекают его главные инженерные преимущества.

\begin{itemize}
    \item \textbf{Простая и быстрая выборка.} Чтобы получить один коэффициент из \(B_\eta\), достаточно извлечь \(2\eta\) случайных битов и выполнить \(\eta\) вычитаний и сложений. Не требуется ни обратной функции ошибок, ни таблиц, ни вычислений с плавающей точкой. В Kyber это реализуется через функцию \texttt{CBD}, которая на вход получает массив байт и выдаёт полином.

    \item \textbf{Выборка за постоянное время.} Алгоритм \texttt{CBD} не содержит ветвлений, зависящих от выбираемых значений, и не использует операций, время выполнения которых менялось бы от значения к значению. Это устраняет тайминг-канал, характерный для наивных реализаций гауссовой выборки.

    \item \textbf{Компактность и детерминированность.} Выборка полностью детерминирована входным массивом байт (который, в свою очередь, порождается псевдослучайной функцией из зерна). Это упрощает анализ и обеспечивает воспроизводимость.

    \item \textbf{Ограниченный носитель.} Коэффициенты CBD всегда лежат в диапазоне \([-\eta,\eta]\), что упрощает анализ вероятности ошибки дешифрования и оценку параметров.
\end{itemize}

\subsection{Почему для Kyber-512 используется \(\eta_1=3\), а для Kyber-768 и Kyber-1024 — \(\eta_1=2\)}

Выбор \(\eta_1\) для секретного вектора \(\mathbf{s}\) и шума \(\mathbf{e}\) (а также для вектора \(\mathbf{r}\) при шифровании) является результатом \textbf{инженерного компромисса} между стойкостью и вероятностью ошибки дешифрования.

\begin{itemize}
    \item \textbf{Kyber-512 (уровень безопасности 1).} Здесь \(k=2\), то есть матрица \(\mathbf{A}\) имеет меньшую размерность, чем в старших наборах. Чтобы обеспечить требуемую стойкость (в терминах сложности решёточных атак, в частности \textit{Core-SVP hardness}) при меньшем \(k\), было решено увеличить разброс секрета и шума, взяв \(\eta_1=3\). Это расширяет носитель до \([-3,3]\) и повышает стандартное отклонение до \(\sqrt{1{,}5}\approx 1{,}225\), что затрудняет атаки на решётку. При этом вероятность ошибки дешифрования остаётся приемлемо низкой (около \(2^{-139}\)).

    \item \textbf{Kyber-768 и Kyber-1024 (уровни безопасности 3 и 5).} Здесь \(k=3\) и \(k=4\) соответственно. Увеличение размерности матрицы \(\mathbf{A}\) само по себе даёт достаточный запас стойкости, поэтому нет необходимости дополнительно расширять распределение шума. Использование \(\eta_1=2\) (носитель \([-2,2]\), \(\sigma\approx 1{,}0\)) позволяет \textbf{снизить вероятность ошибки дешифрования} и одновременно упростить и ускорить выборку. Вероятности ошибки для этих наборов составляют \(2^{-164}\) и \(2^{-174}\) соответственно.
\end{itemize}

Таким образом, \(\eta_1=3\) для Kyber-512 — это способ «добрать» стойкость за счёт несколько большего шума, тогда как для Kyber-768/1024 стойкость обеспечивается размерностью, и \(\eta_1=2\) является более выгодным с точки зрения надёжности дешифрования и производительности.

\subsection{Почему \(\eta_2=2\) для всех уровней}

Параметр \(\eta_2\) управляет распределением вектора \(\mathbf{e}_1\) и скаляра \(e_2\), которые используются только на этапе шифрования (инкапсуляции). Эти шумовые члены не входят в долговременный секретный ключ и не определяют стойкость задачи MLWE в том же смысле, что \(\mathbf{s}\) и \(\mathbf{e}\). Их главная роль — «замаскировать» сообщение и обеспечить семантическую безопасность шифртекста. Поскольку вклад \(\mathbf{e}_1\) и \(e_2\) в общую ошибку дешифрования не является доминирующим, для всех уровней достаточно \(\eta_2=2\). Это упрощает реализацию (не нужно поддерживать разные параметры CBD для шифрования на разных уровнях) и гарантирует низкую вероятность ошибки дешифрования.

\subsection{Сравнение с равномерным распределением: почему не UD}

Иногда в качестве альтернативы CBD рассматривают равномерное распределение (UD) на том же носителе. При одинаковом диапазоне значений UD имеет \textbf{большее стандартное отклонение}, чем CBD (примерно в \(\sqrt{4/3}\) раза). Это могло бы повысить стойкость LWE, но одновременно \textbf{увеличивает вероятность ошибки дешифрования}, что усиливает риск атак, использующих ошибки дешифрования (decryption failure attacks). Для Kyber с его жёсткими требованиями к вероятности ошибки (порядка \(2^{-139}\) и ниже) такой обмен невыгоден. CBD же даёт распределение с \textbf{меньшим стандартным отклонением при том же носителе}, что позволяет сохранить низкую вероятность ошибки при достаточной стойкости.

\subsection{Итог}

Центрированное биномиальное распределение было выбрано в Kyber по трём взаимосвязанным причинам:
\begin{enumerate}
    \item \textbf{Эффективность}: выборка из CBD проста, быстра и не требует высокоточных вычислений или больших таблиц.
    \item \textbf{Безопасность реализации}: CBD можно выбирать за постоянное время, что устраняет тайминг-атаки, характерные для гауссовой выборки.
    \item \textbf{Достаточность для LWE}: стойкость LWE определяется стандартным отклонением, а CBD с нужным \(\eta\) обеспечивает подходящее стандартное отклонение \(\sqrt{\eta/2}\), что подтверждается параметрами и анализом безопасности Kyber.
\end{enumerate}
Различие \(\eta_1=3\) для Kyber-512 и \(\eta_1=2\) для старших уровней отражает настройку этого компромисса под конкретный уровень безопасности, а фиксированное \(\eta_2=2\) для всех уровней связано с вспомогательной ролью шумов шифрования.

\newpage
\section{Rejection sampling для матрицы \(\mathbf{A}\) и \(\mathbf{A}^T\)}

Матрица \(\mathbf{A}\) (и её транспонированная версия \(\mathbf{A}^T\)) генерируется детерминированно из публичного зерна \(\rho\) с помощью rejection sampling (выборки с отбраковкой). В отличие от шумовых векторов, коэффициенты \(\mathbf{A}\) должны быть равномерно распределены в \(\mathbb{Z}_q\) (\(q=3329\)).

\subsection{Пошаговый граф rejection sampling для \(\mathbf{A}\)}

Ниже приведена блок-схема процесса генерации одного полинома матрицы \(\mathbf{A}\) (для одной пары индексов \((i,j)\)). Для полной матрицы процесс повторяется \(k^2\) раз.

\newpage
% \begin{center}
%     \begin{tikzpicture}[
%         node distance = 1.2cm and 1.5cm,
%         box/.style = {rectangle, draw, rounded corners, text width=8cm, align=center, minimum height=0.8cm},
%         decision/.style = {diamond, draw, aspect=2, text width=6cm, align=center, inner sep=1pt},
%         arrow/.style = {-{Stealth[length=2mm]}, thick}
%         ]

%         \node (start) [box] {Вход: зерно \(\rho\), индексы \((i,j)\)};
%         \node (xof) [box, below=of start] {Инициализировать XOF (SHAKE-128) с \(\rho\) и индексами \((i,j)\)};
%         \node (squeeze) [box, below=of xof] {Извлечь 3 байта (24 бита) из потока XOF};
%         \node (split) [box, below=of squeeze] {Разбить 24 бита на два 12-битных кандидата \(c_1, c_2\)};
%         \node (check1) [decision, below=of split] {\(c_1 < 3329\)?};
%         \node (accept1) [box, right=of check1] {Принять \(c_1\) как коэффициент};
%         \node (check2) [decision, below=of check1] {\(c_2 < 3329\)?};
%         \node (accept2) [box, right=of check2] {Принять \(c_2\) как коэффициент};
%         \node (enough) [decision, below=of check2] {Набрано 256 коэффициентов?};
%         \node (output) [box, below=of enough] {Выход: полином \(a_{i,j}(x)\)};

%         \draw [arrow] (start) -- (xof);
%         \draw [arrow] (xof) -- (squeeze);
%         \draw [arrow] (squeeze) -- (split);
%         \draw [arrow] (split) -- (check1);
%         \draw [arrow] (check1) -- node[above] {да} (accept1);
%         \draw [arrow] (check1) -- node[left] {нет} (check2);
%         \draw [arrow] (check2) -- node[above] {да} (accept2);
%         \draw [arrow] (check2) -- node[left] {нет} (enough);
%         \draw [arrow] (accept1) |- (enough);
%         \draw [arrow] (accept2) |- (enough);
%         \draw [arrow] (enough) -- node[left] {нет} (squeeze);
%         \draw [arrow] (enough) -- node[right] {да} (output);

%     \end{tikzpicture}
% \end{center}
\begin{center}
    \begin{tikzpicture}[
        node distance = 1.2cm and 1.5cm,
        box/.style = {rectangle, draw, rounded corners, text width=7cm, align=center, minimum height=0.8cm},
        decision/.style = {diamond, draw, aspect=2, text width=5cm, align=center, inner sep=1pt},
        arrow/.style = {-{Stealth[length=2mm]}, thick},
        lbl/.style = {font=\small, fill=white, inner sep=1pt}
        ]

        \node (start) [box] {Вход: зерно \(\rho\), индексы \((i,j)\)};
        \node (xof) [box, below=of start] {Инициализировать XOF (SHAKE-128) с \(\rho\) и индексами \((i,j)\)};
        \node (squeeze) [box, below=of xof] {Извлечь 3 байта (24 бита) из потока XOF};
        \node (split) [box, below=of squeeze] {Разбить 24 бита на два 12-битных кандидата \(c_1, c_2\)};
        \node (check1) [decision, below=of split] {\(c_1 < 3329\)?};
        \node (check2) [decision, below=of check1] {\(c_2 < 3329\)?};
        \node (enough) [decision, below=of check2] {Набрано 256 коэффициентов?};
        \node (output) [box, below=of enough] {Выход: полином \(a_{i,j}(x)\)};

        \node (accept1) [box, right=of check1] {Принять \(c_1\) как коэффициент};
        \node (accept2) [box, right=of check2] {Принять \(c_2\) как коэффициент};

        % Основной поток сверху вниз
        \draw [arrow] (start) -- (xof);
        \draw [arrow] (xof) -- (squeeze);
        \draw [arrow] (squeeze) -- (split);
        \draw [arrow] (split) -- (check1);

        % check1: да / нет
        \draw [arrow] (check1.east) -- node[lbl, above] {да} (accept1.west);
        \draw [arrow] (check1.south) -- node[lbl, left] {нет} (check2.north);

        % check2: да / нет
        \draw [arrow] (check2.east) -- node[lbl, above] {да} (accept2.west);
        \draw [arrow] (check2.south) -- node[lbl, left] {нет} (enough.north);

        % enough: да вниз
        \draw [arrow] (enough.south) -- node[lbl, right] {да} (output.north);

        % accept1 -> enough: обходим accept2 справа, входим в нижне-правую грань
        \draw [arrow] (accept1.east) -- ++(1.0,0) |- (enough.south east);

        % accept2 -> enough: входим в верхне-правую грань
        \draw [arrow] (accept2.east) -- ++(0.4,0) |- (enough.north east);

        % enough -> squeeze: возврат по левой стороне, минуя все ромбы
        \coordinate (farleft) at ($(enough.west)+(-2.0,0)$);
        \draw [arrow] (enough.west) -- node[lbl, above] {нет} (farleft) |- (squeeze.west);

    \end{tikzpicture}
\end{center}

Вероятность принятия одного 12-битного кандидата:
\[
    P_{\text{принять}} = \frac{3329}{4096} \approx 0{,}8127,
\]
то есть вероятность отбраковки около \(18{,}73\%\).

\subsection{Особенности для транспонированной матрицы \(\mathbf{A}^T\)}

Транспонированная матрица \(\mathbf{A}^T\) генерируется тем же самым алгоритмом rejection sampling, но с \textbf{переставленными индексами}: элемент \(a^T_{i,j}\) соответствует элементу \(a_{j,i}\) исходной матрицы. На практике при генерации \(\mathbf{A}^T\) в качестве входных индексов для XOF используются \((j,i)\) вместо \((i,j)\). Все остальные шаги (извлечение 3 байт, разбиение на два 12-битных кандидата, проверка \(< 3329\), повтор до набора 256 коэффициентов) остаются идентичными.

Таким образом, rejection sampling для \(\mathbf{A}^T\) — это не отдельный алгоритм, а тот же процесс с транспонированной адресацией элементов.

% \section{Заключение}

% В документе приведены точные таблицы вероятностей для всех шумовых векторов Kyber на каждом уровне безопасности, а также детальный граф rejection sampling для матрицы \(\mathbf{A}\) и её транспонированной версии \(\mathbf{A}^T\). Показано, что \(\mathbf{A}^T\) использует тот же механизм отбраковки, но с переставленными индексами.

\end{document}